% --- Archon physics formalization source begin ---
% archon:physics
% archon:covers PhyXMiniProblems/problem_phyx_mini_0948.lean
% archon:source-report reports/phyx_mini/problem_phyx_mini_0948.source.json

\chapter{Physics problem phyx\_mini\_0948}
\label{ch:PhyXMiniProblems_problem_phyx_mini_0948}

\paragraph{Problem source.}
\#\# Physical scenario

The rectangular loop of wire has a mass of \textbackslash{}(0.15\textbackslash{},\textbackslash{}text\{g/cm\}\textbackslash{}) of length and is pivoted about side \textbackslash{}(ab\textbackslash{}) on a frictionless axis. The current in the wire is \textbackslash{}(8.2\textbackslash{},\textbackslash{}text\{A\}\textbackslash{}) in the direction shown.

\#\# Question

Find the magnitude and direction of the magnetic field parallel to the \textbackslash{}(y\textbackslash{})-axis that will cause the loop to swing up until its plane makes an angle of \textbackslash{}(30.0\textasciicircum{}\textbackslash{}circ\textbackslash{}) with the \textbackslash{}(yz\textbackslash{})-plane.

\#\# Answer choices

- A: A: 0.048T

- B: B: 0.024T

- C: C: 0.012T

- D: D: 0.044T

\#\# Figure caption (auxiliary; use the image as primary evidence)

The image shows a three-dimensional geometric figure, specifically a parallelogram, positioned within a coordinate system labeled with the x, y, and z axes. 

Key elements:

1. **Axes**: 
   - The x-axis is horizontal.
   - The y-axis is vertical.
   - The z-axis is diagonal from top left to bottom right in the image.

2. **Parallelogram**:
   - A parallelogram is drawn with a slanted rectangle appearance.
   - The parallelogram is outlined in a golden color.
   - The sides of the parallelogram are not parallel to the axes; they are angled.

3. **Vectors**:
   - A purple vector labeled "a" is directed along the top left side of the parallelogram.
   - A purple vector labeled "b" is directed along the left vertical side of the parallelogram.

4. **Measurements**:
   - One side of the parallelogram has a length labeled as "6.00 cm".
   - Another side has a length labeled as "8.00 cm".

5. **Angle**:
   - The angle between the vector "b" and the z-axis is marked as 30.0°.

This figure seems to represent a vector component problem where the parallelogram rule or law of cosines might be used to solve for vector magnitudes or directions.

\#\# Dataset metadata

Magnetism; Physical Model Grounding Reasoning, Spatial Relation Reasoning

\paragraph{Recorded answer/context.}
B: B: 0.024T

\paragraph{Figure/image path.}
/root/proposal\_for\_physic/science-mango/phyx\_mini\_run/phyx\_data/test\_image/948.png

\paragraph{Formalization target.}
create a compiling Lean file with sorry bodies at `PhyXMiniProblems/problem\_phyx\_mini\_0948.lean`. The Lean declarations must preserve the physical quantities, dimensions or dimensional roles, figure labels, governing-law hypotheses, and final relation expressed by this problem.
Use Mathlib/Physlib names found through LeanExplore where available. If a physics API is missing, introduce faithful local abstractions rather than scalar placeholder aliases.

\begin{theorem}[Physics formalization target]
\label{thm:physics:phyx_mini_0948:target}
\lean{PhyXMiniProblems.ProblemPhyXMini0948.balancing_magnetic_field}
\uses{def:physics:phyx-mini-0948:phyxminiproblems-problemphyxmini0948-pivotedcurrentloopsetup, def:physics:phyx-mini-0948:phyxminiproblems-problemphyxmini0948-matchesprimaryfigure, def:physics:phyx-mini-0948:phyxminiproblems-problemphyxmini0948-matchesproblemdescription, def:physics:phyx-mini-0948:phyxminiproblems-problemphyxmini0948-hasstandardnearearthgravity, def:physics:phyx-mini-0948:phyxminiproblems-problemphyxmini0948-hasphysicalloopparameters, def:physics:phyx-mini-0948:phyxminiproblems-problemphyxmini0948-hasuniformappliedmagneticfield, def:physics:phyx-mini-0948:phyxminiproblems-problemphyxmini0948-satisfiesrectangularwiremasslaw, def:physics:phyx-mini-0948:phyxminiproblems-problemphyxmini0948-satisfiespivottorquemagnitudelaws, def:physics:phyx-mini-0948:phyxminiproblems-problemphyxmini0948-satisfiescurrentloopdirectionlaws, def:physics:phyx-mini-0948:phyxminiproblems-problemphyxmini0948-isinrotationalequilibrium, def:physics:phyx-mini-0948:phyxminiproblems-problemphyxmini0948-magnituderoundstochoice}
The assigned autoformalize agent should translate this physics problem into Lean declarations in the covered file, with theorem and lemma proof bodies written as `by sorry`.
\end{theorem}
\begin{proof}
This is an autoformalization task, not a proof task. Produce statements that can later be proved by the physics prover without weakening the physical model.
\end{proof}

% --- Archon declaration topology begin ---
\section{Declaration topology}
\begin{definition}[magnetic Flux Density Dimension]
  \label{def:physics:phyx-mini-0948:phyxminiproblems-problemphyxmini0948-magneticfluxdensitydimension}
  \lean{PhyXMiniProblems.ProblemPhyXMini0948.magneticFluxDensityDimension}
  The dimension M T⁻¹ C⁻¹ of magnetic flux density (tesla)
\end{definition}
\begin{proof}
  This is the declared model definition of magnetic Flux Density Dimension.
\end{proof}
\begin{definition}[electric Current Dimension]
  \label{def:physics:phyx-mini-0948:phyxminiproblems-problemphyxmini0948-electriccurrentdimension}
  \lean{PhyXMiniProblems.ProblemPhyXMini0948.electricCurrentDimension}
  The dimension C T⁻¹ of electric current (ampere)
\end{definition}
\begin{proof}
  This is the declared model definition of electric Current Dimension.
\end{proof}
\begin{definition}[linear Mass Density Dimension]
  \label{def:physics:phyx-mini-0948:phyxminiproblems-problemphyxmini0948-linearmassdensitydimension}
  \lean{PhyXMiniProblems.ProblemPhyXMini0948.linearMassDensityDimension}
  The dimension M L⁻¹ of linear mass density
\end{definition}
\begin{proof}
  This is the declared model definition of linear Mass Density Dimension.
\end{proof}
\begin{definition}[acceleration Dimension]
  \label{def:physics:phyx-mini-0948:phyxminiproblems-problemphyxmini0948-accelerationdimension}
  \lean{PhyXMiniProblems.ProblemPhyXMini0948.accelerationDimension}
  The dimension L T⁻² of acceleration
\end{definition}
\begin{proof}
  This is the declared model definition of acceleration Dimension.
\end{proof}
\begin{definition}[torque Dimension]
  \label{def:physics:phyx-mini-0948:phyxminiproblems-problemphyxmini0948-torquedimension}
  \lean{PhyXMiniProblems.ProblemPhyXMini0948.torqueDimension}
  The dimension M L² T⁻² of torque
\end{definition}
\begin{proof}
  This is the declared model definition of torque Dimension.
\end{proof}
\begin{definition}[Magnetic Flux Density Magnitude]
  \label{def:physics:phyx-mini-0948:phyxminiproblems-problemphyxmini0948-magneticfluxdensitymagnitude}
  \lean{PhyXMiniProblems.ProblemPhyXMini0948.MagneticFluxDensityMagnitude}
  \uses{def:physics:phyx-mini-0948:phyxminiproblems-problemphyxmini0948-magneticfluxdensitydimension}
  A nonnegative, unit-independent magnetic-flux-density magnitude
\end{definition}
\begin{proof}
  This is the declared model definition of Magnetic Flux Density Magnitude.
\end{proof}
\begin{definition}[Electric Current Magnitude]
  \label{def:physics:phyx-mini-0948:phyxminiproblems-problemphyxmini0948-electriccurrentmagnitude}
  \lean{PhyXMiniProblems.ProblemPhyXMini0948.ElectricCurrentMagnitude}
  \uses{def:physics:phyx-mini-0948:phyxminiproblems-problemphyxmini0948-electriccurrentdimension}
  A nonnegative, unit-independent electric-current magnitude
\end{definition}
\begin{proof}
  This is the declared model definition of Electric Current Magnitude.
\end{proof}
\begin{definition}[Length Magnitude]
  \label{def:physics:phyx-mini-0948:phyxminiproblems-problemphyxmini0948-lengthmagnitude}
  \lean{PhyXMiniProblems.ProblemPhyXMini0948.LengthMagnitude}
  A nonnegative, unit-independent physical length
\end{definition}
\begin{proof}
  This is the declared model definition of Length Magnitude.
\end{proof}
\begin{definition}[Mass Magnitude]
  \label{def:physics:phyx-mini-0948:phyxminiproblems-problemphyxmini0948-massmagnitude}
  \lean{PhyXMiniProblems.ProblemPhyXMini0948.MassMagnitude}
  A nonnegative, unit-independent mass
\end{definition}
\begin{proof}
  This is the declared model definition of Mass Magnitude.
\end{proof}
\begin{definition}[Linear Mass Density Magnitude]
  \label{def:physics:phyx-mini-0948:phyxminiproblems-problemphyxmini0948-linearmassdensitymagnitude}
  \lean{PhyXMiniProblems.ProblemPhyXMini0948.LinearMassDensityMagnitude}
  \uses{def:physics:phyx-mini-0948:phyxminiproblems-problemphyxmini0948-linearmassdensitydimension}
  A nonnegative, unit-independent mass per unit length
\end{definition}
\begin{proof}
  This is the declared model definition of Linear Mass Density Magnitude.
\end{proof}
\begin{definition}[Acceleration Magnitude]
  \label{def:physics:phyx-mini-0948:phyxminiproblems-problemphyxmini0948-accelerationmagnitude}
  \lean{PhyXMiniProblems.ProblemPhyXMini0948.AccelerationMagnitude}
  \uses{def:physics:phyx-mini-0948:phyxminiproblems-problemphyxmini0948-accelerationdimension}
  A nonnegative, unit-independent acceleration magnitude
\end{definition}
\begin{proof}
  This is the declared model definition of Acceleration Magnitude.
\end{proof}
\begin{definition}[Torque Magnitude]
  \label{def:physics:phyx-mini-0948:phyxminiproblems-problemphyxmini0948-torquemagnitude}
  \lean{PhyXMiniProblems.ProblemPhyXMini0948.TorqueMagnitude}
  \uses{def:physics:phyx-mini-0948:phyxminiproblems-problemphyxmini0948-torquedimension}
  A nonnegative, unit-independent torque magnitude
\end{definition}
\begin{proof}
  This is the declared model definition of Torque Magnitude.
\end{proof}
\begin{definition}[magnetic Flux Density In Teslas]
  \label{def:physics:phyx-mini-0948:phyxminiproblems-problemphyxmini0948-magneticfluxdensityinteslas}
  \lean{PhyXMiniProblems.ProblemPhyXMini0948.magneticFluxDensityInTeslas}
  \uses{def:physics:phyx-mini-0948:phyxminiproblems-problemphyxmini0948-magneticfluxdensitymagnitude}
  Coherent-SI readout of magnetic flux density, in teslas
\end{definition}
\begin{proof}
  This is the declared model definition of magnetic Flux Density In Teslas.
\end{proof}
\begin{definition}[electric Current In Amperes]
  \label{def:physics:phyx-mini-0948:phyxminiproblems-problemphyxmini0948-electriccurrentinamperes}
  \lean{PhyXMiniProblems.ProblemPhyXMini0948.electricCurrentInAmperes}
  \uses{def:physics:phyx-mini-0948:phyxminiproblems-problemphyxmini0948-electriccurrentmagnitude}
  Coherent-SI readout of electric current, in amperes
\end{definition}
\begin{proof}
  This is the declared model definition of electric Current In Amperes.
\end{proof}
\begin{definition}[length In Meters]
  \label{def:physics:phyx-mini-0948:phyxminiproblems-problemphyxmini0948-lengthinmeters}
  \lean{PhyXMiniProblems.ProblemPhyXMini0948.lengthInMeters}
  \uses{def:physics:phyx-mini-0948:phyxminiproblems-problemphyxmini0948-lengthmagnitude}
  Coherent-SI readout of length, in metres
\end{definition}
\begin{proof}
  This is the declared model definition of length In Meters.
\end{proof}
\begin{definition}[length In Centimeters]
  \label{def:physics:phyx-mini-0948:phyxminiproblems-problemphyxmini0948-lengthincentimeters}
  \lean{PhyXMiniProblems.ProblemPhyXMini0948.lengthInCentimeters}
  \uses{def:physics:phyx-mini-0948:phyxminiproblems-problemphyxmini0948-lengthmagnitude, def:physics:phyx-mini-0948:phyxminiproblems-problemphyxmini0948-lengthinmeters}
  Centimetre readout obtained from the coherent-SI length readout
\end{definition}
\begin{proof}
  This is the declared model definition of length In Centimeters.
\end{proof}
\begin{definition}[mass In Kilograms]
  \label{def:physics:phyx-mini-0948:phyxminiproblems-problemphyxmini0948-massinkilograms}
  \lean{PhyXMiniProblems.ProblemPhyXMini0948.massInKilograms}
  \uses{def:physics:phyx-mini-0948:phyxminiproblems-problemphyxmini0948-massmagnitude}
  Coherent-SI readout of mass, in kilograms
\end{definition}
\begin{proof}
  This is the declared model definition of mass In Kilograms.
\end{proof}
\begin{definition}[linear Mass Density In Kilograms Per Meter]
  \label{def:physics:phyx-mini-0948:phyxminiproblems-problemphyxmini0948-linearmassdensityinkilogramspermeter}
  \lean{PhyXMiniProblems.ProblemPhyXMini0948.linearMassDensityInKilogramsPerMeter}
  \uses{def:physics:phyx-mini-0948:phyxminiproblems-problemphyxmini0948-linearmassdensitymagnitude}
  Coherent-SI readout of linear mass density, in kilograms per metre
\end{definition}
\begin{proof}
  This is the declared model definition of linear Mass Density In Kilograms Per Meter.
\end{proof}
\begin{definition}[linear Mass Density In Grams Per Centimeter]
  \label{def:physics:phyx-mini-0948:phyxminiproblems-problemphyxmini0948-linearmassdensityingramspercentimeter}
  \lean{PhyXMiniProblems.ProblemPhyXMini0948.linearMassDensityInGramsPerCentimeter}
  \uses{def:physics:phyx-mini-0948:phyxminiproblems-problemphyxmini0948-linearmassdensitymagnitude, def:physics:phyx-mini-0948:phyxminiproblems-problemphyxmini0948-linearmassdensityinkilogramspermeter}
  Grams-per-centimetre readout; 1 kg/m = 10 g/cm
\end{definition}
\begin{proof}
  This is the declared model definition of linear Mass Density In Grams Per Centimeter.
\end{proof}
\begin{definition}[acceleration In Meters Per Second Squared]
  \label{def:physics:phyx-mini-0948:phyxminiproblems-problemphyxmini0948-accelerationinmeterspersecondsquared}
  \lean{PhyXMiniProblems.ProblemPhyXMini0948.accelerationInMetersPerSecondSquared}
  \uses{def:physics:phyx-mini-0948:phyxminiproblems-problemphyxmini0948-accelerationmagnitude}
  Coherent-SI readout of acceleration, in metres per second squared
\end{definition}
\begin{proof}
  This is the declared model definition of acceleration In Meters Per Second Squared.
\end{proof}
\begin{definition}[torque In Newton Meters]
  \label{def:physics:phyx-mini-0948:phyxminiproblems-problemphyxmini0948-torqueinnewtonmeters}
  \lean{PhyXMiniProblems.ProblemPhyXMini0948.torqueInNewtonMeters}
  \uses{def:physics:phyx-mini-0948:phyxminiproblems-problemphyxmini0948-torquemagnitude}
  Coherent-SI readout of torque, in newton metres
\end{definition}
\begin{proof}
  This is the declared model definition of torque In Newton Meters.
\end{proof}
\begin{definition}[Coordinate Axis]
  \label{def:physics:phyx-mini-0948:phyxminiproblems-problemphyxmini0948-coordinateaxis}
  \lean{PhyXMiniProblems.ProblemPhyXMini0948.CoordinateAxis}
  The coordinate axes printed in the primary figure
\end{definition}
\begin{proof}
  This is the declared model definition of Coordinate Axis.
\end{proof}
\begin{definition}[Signed Axis Direction]
  \label{def:physics:phyx-mini-0948:phyxminiproblems-problemphyxmini0948-signedaxisdirection}
  \lean{PhyXMiniProblems.ProblemPhyXMini0948.SignedAxisDirection}
  The six signed coordinate-axis directions
\end{definition}
\begin{proof}
  This is the declared model definition of Signed Axis Direction.
\end{proof}
\begin{definition}[axis]
  \label{def:physics:phyx-mini-0948:phyxminiproblems-problemphyxmini0948-signedaxisdirection-axis}
  \lean{PhyXMiniProblems.ProblemPhyXMini0948.SignedAxisDirection.axis}
  \uses{def:physics:phyx-mini-0948:phyxminiproblems-problemphyxmini0948-coordinateaxis, def:physics:phyx-mini-0948:phyxminiproblems-problemphyxmini0948-signedaxisdirection}
  The unsigned coordinate axis underlying a signed direction
\end{definition}
\begin{proof}
  This is the declared model definition of axis.
\end{proof}
\begin{definition}[Spatial Vector]
  \label{def:physics:phyx-mini-0948:phyxminiproblems-problemphyxmini0948-spatialvector}
  \lean{PhyXMiniProblems.ProblemPhyXMini0948.SpatialVector}
  Cartesian spatial vectors used for the field and right-hand rules
\end{definition}
\begin{proof}
  This is the declared model definition of Spatial Vector.
\end{proof}
\begin{definition}[unit Vector]
  \label{def:physics:phyx-mini-0948:phyxminiproblems-problemphyxmini0948-signedaxisdirection-unitvector}
  \lean{PhyXMiniProblems.ProblemPhyXMini0948.SignedAxisDirection.unitVector}
  \uses{def:physics:phyx-mini-0948:phyxminiproblems-problemphyxmini0948-signedaxisdirection, def:physics:phyx-mini-0948:phyxminiproblems-problemphyxmini0948-spatialvector}
  Unit vector along a signed coordinate-axis direction
\end{definition}
\begin{proof}
  This is the declared model definition of unit Vector.
\end{proof}
\begin{definition}[Has Signed Axis Direction]
  \label{def:physics:phyx-mini-0948:phyxminiproblems-problemphyxmini0948-hassignedaxisdirection}
  \lean{PhyXMiniProblems.ProblemPhyXMini0948.HasSignedAxisDirection}
  \uses{def:physics:phyx-mini-0948:phyxminiproblems-problemphyxmini0948-signedaxisdirection, def:physics:phyx-mini-0948:phyxminiproblems-problemphyxmini0948-spatialvector}
  A nonzero vector is a positive multiple of the indicated signed axis
\end{definition}
\begin{proof}
  This is the declared model definition of Has Signed Axis Direction.
\end{proof}
\begin{definition}[Crosses To Direction]
  \label{def:physics:phyx-mini-0948:phyxminiproblems-problemphyxmini0948-crossestodirection}
  \lean{PhyXMiniProblems.ProblemPhyXMini0948.CrossesToDirection}
  \uses{def:physics:phyx-mini-0948:phyxminiproblems-problemphyxmini0948-signedaxisdirection}
  The Physlib cross product of two signed unit directions has a third one
\end{definition}
\begin{proof}
  This is the declared model definition of Crosses To Direction.
\end{proof}
\begin{definition}[Loop Vertex]
  \label{def:physics:phyx-mini-0948:phyxminiproblems-problemphyxmini0948-loopvertex}
  \lean{PhyXMiniProblems.ProblemPhyXMini0948.LoopVertex}
  Lettered vertices of the rectangular loop
\end{definition}
\begin{proof}
  This is the declared model definition of Loop Vertex.
\end{proof}
\begin{definition}[Loop Edge]
  \label{def:physics:phyx-mini-0948:phyxminiproblems-problemphyxmini0948-loopedge}
  \lean{PhyXMiniProblems.ProblemPhyXMini0948.LoopEdge}
  Oriented boundary edges, named by their first and second endpoint
\end{definition}
\begin{proof}
  This is the declared model definition of Loop Edge.
\end{proof}
\begin{definition}[Edge Traversal]
  \label{def:physics:phyx-mini-0948:phyxminiproblems-problemphyxmini0948-edgetraversal}
  \lean{PhyXMiniProblems.ProblemPhyXMini0948.EdgeTraversal}
  Which way current traverses a named boundary edge
\end{definition}
\begin{proof}
  This is the declared model definition of Edge Traversal.
\end{proof}
\begin{definition}[Loop Shape]
  \label{def:physics:phyx-mini-0948:phyxminiproblems-problemphyxmini0948-loopshape}
  \lean{PhyXMiniProblems.ProblemPhyXMini0948.LoopShape}
  Shape classification of the current-carrying wire loop
\end{definition}
\begin{proof}
  This is the declared model definition of Loop Shape.
\end{proof}
\begin{definition}[Rectangular Loop Figure]
  \label{def:physics:phyx-mini-0948:phyxminiproblems-problemphyxmini0948-rectangularloopfigure}
  \lean{PhyXMiniProblems.ProblemPhyXMini0948.RectangularLoopFigure}
  \uses{def:physics:phyx-mini-0948:phyxminiproblems-problemphyxmini0948-coordinateaxis, def:physics:phyx-mini-0948:phyxminiproblems-problemphyxmini0948-signedaxisdirection, def:physics:phyx-mini-0948:phyxminiproblems-problemphyxmini0948-loopedge, def:physics:phyx-mini-0948:phyxminiproblems-problemphyxmini0948-edgetraversal}
  Literal geometric and directional information carried by image 948.png. The unknown magnetic field is deliberately absent from this record
\end{definition}
\begin{proof}
  This is the declared model definition of Rectangular Loop Figure.
\end{proof}
\begin{definition}[Pivoted Current Loop Setup]
  \label{def:physics:phyx-mini-0948:phyxminiproblems-problemphyxmini0948-pivotedcurrentloopsetup}
  \lean{PhyXMiniProblems.ProblemPhyXMini0948.PivotedCurrentLoopSetup}
  \uses{def:physics:phyx-mini-0948:phyxminiproblems-problemphyxmini0948-magneticfluxdensitymagnitude, def:physics:phyx-mini-0948:phyxminiproblems-problemphyxmini0948-electriccurrentmagnitude, def:physics:phyx-mini-0948:phyxminiproblems-problemphyxmini0948-lengthmagnitude, def:physics:phyx-mini-0948:phyxminiproblems-problemphyxmini0948-massmagnitude, def:physics:phyx-mini-0948:phyxminiproblems-problemphyxmini0948-linearmassdensitymagnitude, def:physics:phyx-mini-0948:phyxminiproblems-problemphyxmini0948-accelerationmagnitude, def:physics:phyx-mini-0948:phyxminiproblems-problemphyxmini0948-torquemagnitude, def:physics:phyx-mini-0948:phyxminiproblems-problemphyxmini0948-signedaxisdirection, def:physics:phyx-mini-0948:phyxminiproblems-problemphyxmini0948-loopshape, def:physics:phyx-mini-0948:phyxminiproblems-problemphyxmini0948-rectangularloopfigure}
  Independent physical quantities describing the pivoted current loop
\end{definition}
\begin{proof}
  This is the declared model definition of Pivoted Current Loop Setup.
\end{proof}
\begin{definition}[Matches Primary Figure]
  \label{def:physics:phyx-mini-0948:phyxminiproblems-problemphyxmini0948-matchesprimaryfigure}
  \lean{PhyXMiniProblems.ProblemPhyXMini0948.MatchesPrimaryFigure}
  \uses{def:physics:phyx-mini-0948:phyxminiproblems-problemphyxmini0948-pivotedcurrentloopsetup}
  Primary-raster evidence: axes, edge labels, current arrows, the 30.0° swing, and the direction of that swing. It contains no magnetic-field answer
\end{definition}
\begin{proof}
  This is the declared model definition of Matches Primary Figure.
\end{proof}
\begin{definition}[Matches Problem Description]
  \label{def:physics:phyx-mini-0948:phyxminiproblems-problemphyxmini0948-matchesproblemdescription}
  \lean{PhyXMiniProblems.ProblemPhyXMini0948.MatchesProblemDescription}
  \uses{def:physics:phyx-mini-0948:phyxminiproblems-problemphyxmini0948-electriccurrentinamperes, def:physics:phyx-mini-0948:phyxminiproblems-problemphyxmini0948-lengthincentimeters, def:physics:phyx-mini-0948:phyxminiproblems-problemphyxmini0948-linearmassdensityingramspercentimeter, def:physics:phyx-mini-0948:phyxminiproblems-problemphyxmini0948-pivotedcurrentloopsetup}
  Problem-statement data and conversions. fieldIsParallelToYAxis restricts the unknown field to one of the two y directions but does not choose its sign
\end{definition}
\begin{proof}
  This is the declared model definition of Matches Problem Description.
\end{proof}
\begin{definition}[Has Standard Near Earth Gravity]
  \label{def:physics:phyx-mini-0948:phyxminiproblems-problemphyxmini0948-hasstandardnearearthgravity}
  \lean{PhyXMiniProblems.ProblemPhyXMini0948.HasStandardNearEarthGravity}
  \uses{def:physics:phyx-mini-0948:phyxminiproblems-problemphyxmini0948-accelerationinmeterspersecondsquared, def:physics:phyx-mini-0948:phyxminiproblems-problemphyxmini0948-pivotedcurrentloopsetup}
  Standard near-Earth gravitational acceleration used by the textbook model
\end{definition}
\begin{proof}
  This is the declared model definition of Has Standard Near Earth Gravity.
\end{proof}
\begin{definition}[Has Physical Loop Parameters]
  \label{def:physics:phyx-mini-0948:phyxminiproblems-problemphyxmini0948-hasphysicalloopparameters}
  \lean{PhyXMiniProblems.ProblemPhyXMini0948.HasPhysicalLoopParameters}
  \uses{def:physics:phyx-mini-0948:phyxminiproblems-problemphyxmini0948-magneticfluxdensityinteslas, def:physics:phyx-mini-0948:phyxminiproblems-problemphyxmini0948-electriccurrentinamperes, def:physics:phyx-mini-0948:phyxminiproblems-problemphyxmini0948-lengthinmeters, def:physics:phyx-mini-0948:phyxminiproblems-problemphyxmini0948-massinkilograms, def:physics:phyx-mini-0948:phyxminiproblems-problemphyxmini0948-linearmassdensityinkilogramspermeter, def:physics:phyx-mini-0948:phyxminiproblems-problemphyxmini0948-accelerationinmeterspersecondsquared, def:physics:phyx-mini-0948:phyxminiproblems-problemphyxmini0948-pivotedcurrentloopsetup}
  Positivity and nondegeneracy conditions for the physical equilibrium
\end{definition}
\begin{proof}
  This is the declared model definition of Has Physical Loop Parameters.
\end{proof}
\begin{definition}[Has Uniform Applied Magnetic Field]
  \label{def:physics:phyx-mini-0948:phyxminiproblems-problemphyxmini0948-hasuniformappliedmagneticfield}
  \lean{PhyXMiniProblems.ProblemPhyXMini0948.HasUniformAppliedMagneticField}
  \uses{def:physics:phyx-mini-0948:phyxminiproblems-problemphyxmini0948-magneticfluxdensityinteslas, def:physics:phyx-mini-0948:phyxminiproblems-problemphyxmini0948-hassignedaxisdirection, def:physics:phyx-mini-0948:phyxminiproblems-problemphyxmini0948-pivotedcurrentloopsetup}
  The Physlib magnetic vector field is uniform in magnitude and has the signed axis direction recorded independently in the setup
\end{definition}
\begin{proof}
  This is the declared model definition of Has Uniform Applied Magnetic Field.
\end{proof}
\begin{definition}[Satisfies Rectangular Wire Mass Law]
  \label{def:physics:phyx-mini-0948:phyxminiproblems-problemphyxmini0948-satisfiesrectangularwiremasslaw}
  \lean{PhyXMiniProblems.ProblemPhyXMini0948.SatisfiesRectangularWireMassLaw}
  \uses{def:physics:phyx-mini-0948:phyxminiproblems-problemphyxmini0948-lengthinmeters, def:physics:phyx-mini-0948:phyxminiproblems-problemphyxmini0948-massinkilograms, def:physics:phyx-mini-0948:phyxminiproblems-problemphyxmini0948-linearmassdensityinkilogramspermeter, def:physics:phyx-mini-0948:phyxminiproblems-problemphyxmini0948-pivotedcurrentloopsetup}
  The mass of the wire is its linear density times its rectangular perimeter
\end{definition}
\begin{proof}
  This is the declared model definition of Satisfies Rectangular Wire Mass Law.
\end{proof}
\begin{definition}[Satisfies Pivot Torque Magnitude Laws]
  \label{def:physics:phyx-mini-0948:phyxminiproblems-problemphyxmini0948-satisfiespivottorquemagnitudelaws}
  \lean{PhyXMiniProblems.ProblemPhyXMini0948.SatisfiesPivotTorqueMagnitudeLaws}
  \uses{def:physics:phyx-mini-0948:phyxminiproblems-problemphyxmini0948-magneticfluxdensityinteslas, def:physics:phyx-mini-0948:phyxminiproblems-problemphyxmini0948-electriccurrentinamperes, def:physics:phyx-mini-0948:phyxminiproblems-problemphyxmini0948-lengthinmeters, def:physics:phyx-mini-0948:phyxminiproblems-problemphyxmini0948-massinkilograms, def:physics:phyx-mini-0948:phyxminiproblems-problemphyxmini0948-accelerationinmeterspersecondsquared, def:physics:phyx-mini-0948:phyxminiproblems-problemphyxmini0948-torqueinnewtonmeters, def:physics:phyx-mini-0948:phyxminiproblems-problemphyxmini0948-pivotedcurrentloopsetup}
  Torque-magnitude laws about the frictionless pivot. The angle is measured from the yz reference plane. The magnetic moment's component perpendicular to a y-directed field contributes cos θ; the center-of-mass lever arm for gravity contributes sin θ
\end{definition}
\begin{proof}
  This is the declared model definition of Satisfies Pivot Torque Magnitude Laws.
\end{proof}
\begin{definition}[Satisfies Current Loop Direction Laws]
  \label{def:physics:phyx-mini-0948:phyxminiproblems-problemphyxmini0948-satisfiescurrentloopdirectionlaws}
  \lean{PhyXMiniProblems.ProblemPhyXMini0948.SatisfiesCurrentLoopDirectionLaws}
  \uses{def:physics:phyx-mini-0948:phyxminiproblems-problemphyxmini0948-crossestodirection, def:physics:phyx-mini-0948:phyxminiproblems-problemphyxmini0948-pivotedcurrentloopsetup}
  General right-hand rules for this oriented loop. The first equation obtains the reference magnetic-moment normal from the two current-oriented sides. The second is the direction law τ = μ × B; the component of μ parallel to B at nonzero tilt contributes no torque. The third relates torque sign to the shown swing
\end{definition}
\begin{proof}
  This is the declared model definition of Satisfies Current Loop Direction Laws.
\end{proof}
\begin{definition}[Is In Rotational Equilibrium]
  \label{def:physics:phyx-mini-0948:phyxminiproblems-problemphyxmini0948-isinrotationalequilibrium}
  \lean{PhyXMiniProblems.ProblemPhyXMini0948.IsInRotationalEquilibrium}
  \uses{def:physics:phyx-mini-0948:phyxminiproblems-problemphyxmini0948-torqueinnewtonmeters, def:physics:phyx-mini-0948:phyxminiproblems-problemphyxmini0948-pivotedcurrentloopsetup}
  Static rotational equilibrium about the frictionless side ab
\end{definition}
\begin{proof}
  This is the declared model definition of Is In Rotational Equilibrium.
\end{proof}
\begin{lemma}[equilibrium field magnitude formula]
  \label{lem:physics:phyx-mini-0948:phyxminiproblems-problemphyxmini0948-equilibrium-field-magnitude-formula}
  \lean{PhyXMiniProblems.ProblemPhyXMini0948.equilibrium_field_magnitude_formula}
  \uses{def:physics:phyx-mini-0948:phyxminiproblems-problemphyxmini0948-magneticfluxdensityinteslas, def:physics:phyx-mini-0948:phyxminiproblems-problemphyxmini0948-electriccurrentinamperes, def:physics:phyx-mini-0948:phyxminiproblems-problemphyxmini0948-lengthinmeters, def:physics:phyx-mini-0948:phyxminiproblems-problemphyxmini0948-linearmassdensityinkilogramspermeter, def:physics:phyx-mini-0948:phyxminiproblems-problemphyxmini0948-accelerationinmeterspersecondsquared, def:physics:phyx-mini-0948:phyxminiproblems-problemphyxmini0948-pivotedcurrentloopsetup, def:physics:phyx-mini-0948:phyxminiproblems-problemphyxmini0948-hasphysicalloopparameters, def:physics:phyx-mini-0948:phyxminiproblems-problemphyxmini0948-satisfiesrectangularwiremasslaw, def:physics:phyx-mini-0948:phyxminiproblems-problemphyxmini0948-satisfiespivottorquemagnitudelaws, def:physics:phyx-mini-0948:phyxminiproblems-problemphyxmini0948-isinrotationalequilibrium}
  The general scalar consequence of wire mass, the two torque laws, and torque balance. It is a derived relation, not a premise containing the requested numeric answer
\end{lemma}
\begin{proof}
  The conclusion follows from the governing relations and figure readouts named in the statement of equilibrium field magnitude formula.
\end{proof}
\begin{definition}[Answer Choice]
  \label{def:physics:phyx-mini-0948:phyxminiproblems-problemphyxmini0948-answerchoice}
  \lean{PhyXMiniProblems.ProblemPhyXMini0948.AnswerChoice}
  Labels of the four displayed magnetic-field choices
\end{definition}
\begin{proof}
  This is the declared model definition of Answer Choice.
\end{proof}
\begin{definition}[field Magnitude In Teslas]
  \label{def:physics:phyx-mini-0948:phyxminiproblems-problemphyxmini0948-answerchoice-fieldmagnitudeinteslas}
  \lean{PhyXMiniProblems.ProblemPhyXMini0948.AnswerChoice.fieldMagnitudeInTeslas}
  \uses{def:physics:phyx-mini-0948:phyxminiproblems-problemphyxmini0948-answerchoice}
  Magnetic-flux-density magnitude in teslas printed by each answer choice
\end{definition}
\begin{proof}
  This is the declared model definition of field Magnitude In Teslas.
\end{proof}
\begin{definition}[Magnitude Rounds To Choice]
  \label{def:physics:phyx-mini-0948:phyxminiproblems-problemphyxmini0948-magnituderoundstochoice}
  \lean{PhyXMiniProblems.ProblemPhyXMini0948.MagnitudeRoundsToChoice}
  \uses{def:physics:phyx-mini-0948:phyxminiproblems-problemphyxmini0948-magneticfluxdensityinteslas, def:physics:phyx-mini-0948:phyxminiproblems-problemphyxmini0948-pivotedcurrentloopsetup, def:physics:phyx-mini-0948:phyxminiproblems-problemphyxmini0948-answerchoice}
  A computed field agrees with a three-decimal-place answer choice. The strict half-millitesla tolerance is the usual rounding interval for the displayed precision
\end{definition}
\begin{proof}
  This is the declared model definition of Magnitude Rounds To Choice.
\end{proof}
% --- Archon declaration topology end ---
% --- Archon physics formalization source end ---
